% No preamble. Prefix olc-. Fixed-regulator leading stationary amplitudes.
\section{The actual leading determinant correction satisfies a cocycle}
\label{olc-section}

This section proves an exact identity of normalized leading stationary
amplitudes at a fixed finite regulator. It compares these amplitudes
through one auxiliary local orbit-tube integral. It does not identify
direct and sequential sharp-window sectors at positive coupling.
All maps, pivot factors and retained Hessians are those of the
actual composite ordinary constraint.

\subsection{Objects and statement}
\label{olc-setup-section}

Fix a finite compatible periodic regulator, a preceding homogeneous
depth $M$, and two allowed ratios $R=\ell^r$, $L=\ell^s$.
Let $X,Y,Z$ be the current, $R$-coarse and $RL$-coarse link
configuration spaces, each with product normalized Haar measure.
The actual iterated ordinary maps satisfy
\begin{equation}
 f=F_R:X\longrightarrow Y,\qquad
 h=F_L:Y\longrightarrow Z,\qquad
 h\circ f=F_{RL}.
 \label{olc-maps}
\end{equation}
No equality with a direct-$R$ or direct-$RL$ ordinary average
is used.

The coupling-free action $S_M$ is the established homogeneous
ordinary minimum action, normalized at the identity. For an
allowed ratio $K$, write $V_*^{M,K}(W)$ for the physical
gauge-zero stationary representative, using the independently
chosen direct-$K$ internal contour gauge $P_K$. Let
$A_{M,K}(W)$ be the full retained Hessian of $S_M+G_K$ there,
in the actual constraint chart. Let $d_K(V_*)$ be the product
of inverse pivot Jacobians. Define
\begin{equation}
 D_{M,K}(W)
  =\log\frac{d_K(V_*^{M,K}(W))}{d_K(1)}
   -\frac12\log\det\!\left[
         A_{M,K}(W)A_{M,K}(0)^{-1}\right].
 \label{olc-leading}
\end{equation}
Here $D_{M,K}(1)=0$. These are the actual $\Delta+\mathcal Q_{\det}$
terms, with positive real branches near the identity and their
given analytic continuations. Constraint curvature and the
absolute contour-gauge Hessian are included in $A_{M,K}$.

Let $V_p^{MR,L}(W)$ be the outer pre-root-gauge stationary
representative, and let $\mathcal T_{MR,L}$ denote substitution
into this field.

\begin{theorem}[Finite leading-amplitude cocycle]
\label{olc-theorem}
On a real small-field germ,
\begin{equation}
 D_{M,RL}
   =\mathcal T_{MR,L}D_{M,R}+D_{MR,L}.
 \label{olc-cocycle}
\end{equation}
The identity extends to every connected common analytic domain
already available for its two sides. The germ used in its proof
may depend on the fixed regulator and the ratios.
\end{theorem}

The proof keeps all Haar, gauge and Gaussian normalization
factors. They disappear only when each leading amplitude is
divided by its value at the identity.

\subsection{Based gauge groups and an exact insertion}
\label{olc-groups-section}

Let $H_R$ consist of current site gauges equal to the identity
at every intermediate $R$-center, and let $H_{RL}$ be based
at the final $RL$-centers. Restriction to intermediate centers
gives
\begin{equation}
 1\longrightarrow H_R\longrightarrow H_{RL}
       \longrightarrow H_L\longrightarrow1,
 \label{olc-group-sequence}
\end{equation}
where $H_L$ is the intermediate group based at final centers.
The site coordinates split into the non-intermediate sites
and the intermediate non-final centers. Embed $H_L$ by setting
its other current-site values equal to the identity. The two
factors have disjoint supports, so their direct product is
$H_{RL}$ and normalized Haar is their product Haar.
There is no background-dependent orbit-volume factor.

Each based action is free. A gauge preserving every bond of a
connected graph is determined by its value at one site; its
value at a based center is the identity.
The actual maps obey
\begin{equation}
 f(U^u)=f(U)^{u|_{\text{intermediate centers}}},\qquad
 h(V^v)=h(V)^{v|_{\text{final centers}}}.
 \label{olc-equivariance}
\end{equation}
Thus $f$ is $H_R$-invariant, $h$ is $H_L$-invariant, and
$h\circ f$ is $H_{RL}$-invariant. Product link Haar is
invariant under all these site actions.

For $u\in H_K$, endpoint covariance of the actual direct-$K$
contour variables gives
\[
 P_K(U^u;y,x)=P_K(U;y,x)u(x)^{-1}.
\]
With the actual nonnegative class penalty
$q(v)=1-\operatorname{Re}\tau(v)$, whose unique zero is the
identity, put $G_K(U)=\sum_{x\ {\rm not\ a\ center}}q(P_K(U;y,x))$.
Then
\begin{equation}
 \int_{H_K}e^{-G_K(U^u)/g^2}\,du
       =c(g)^{n_{s,K}},\qquad
 c(g)=\int_G e^{-q(v)/g^2}\,dv .
 \label{olc-gauge-insertion}
\end{equation}
It is therefore exact to insert $e^{-G_K/g^2}$ into an
$H_K$-invariant tube integral and multiply by
$c(g)^{-n_{s,K}}$. This factor is independent of the target
field. Its asymptotic constant and power of $g$ need not be
computed below.

\subsection{A consistent local orbit-tube phase}
\label{olc-tube-section}

The original $S_M$ is known on a small-link domain. We first
justify extending it to the particular saturated tube needed
for an exact pushforward composition.

Suppose two sufficiently small configurations $U,V$ satisfy
$V=U^u$ with $u\in H_{RL}$. The bond equation, in a
unitary-invariant matrix norm, yields
\begin{equation}
 \|u(x)-u(x+e_\mu)\|
   \leq C\bigl(\|U(x,\mu)-I\|+\|V(x,\mu)-I\|\bigr).
 \label{olc-gauge-edge}
\end{equation}
Every site has a path of at most $C_dRL$ edges to a final
center, where $u=I$. Consequently
\begin{equation}
 \|u-I\|_\infty
   \leq C_dRL\bigl(\|U-I\|_\infty+\|V-I\|_\infty\bigr).
 \label{olc-gauge-small}
\end{equation}
Shrink the comparison ball so that the near-identity logarithm
$\lambda=\log u$ exists. Its edge differences satisfy the
same type of estimate as \eqref{olc-gauge-edge}.
Along $u_t=\exp(t\lambda)$ the transformed links stay within
a fixed multiple of the endpoint smallness. Fixed color-norm
conversions are included in this constant. Shrinking the ball
again keeps the whole path in the established action domain.
Local Ward invariance therefore gives
\begin{equation}
                         S_M(U)=S_M(V).
 \label{olc-phase-consistency}
\end{equation}
The same argument applies on the intermediate graph and to
$H_L$. No radius uniform in $RL$ is claimed in this step.

On the $H_{RL}$-saturation of this small representative
neighborhood, define
\begin{equation}
                       S_{\rm ext}(U^u)=S_M(U).
 \label{olc-extended-phase}
\end{equation}
Equation \eqref{olc-phase-consistency} makes the definition
independent of the small representative. The free compact
group action has an actual local slice $P_{RL}=I$: the unique
based transformation to it is $u(x)=P_{RL}(U;y,x)$.
Hence $S_{\rm ext}$ is smooth, and real analytic in local
slice coordinates, and is invariant on the saturation.

The local ordinary maps and contour variables extend
equivariantly to this tube. One may use their endpoint-covariant
relative-log formulas or define them from the slice. On
overlaps the two definitions agree by
\eqref{olc-equivariance} and the original local formulas.
No measurable completion away from the tube is invoked.

The proved physical branch is a strict local constrained
minimum modulo the based orbit. Its actual retained
Hessian after adding $G_K$ is positive. On the slice $P_K=I$,
the gauge penalty is identically zero, so that positivity
also gives positivity of the action Hessian on the quotient
constrained tangent. The finite-dimensional implicit-function
argument yields a unique local critical orbit in a smaller
neighborhood.

Choose a compact closure of a sufficiently small real target
germ. Its stationary orbits form a compact set. Choose a smooth
$H_{RL}$-invariant cutoff $\chi$, equal to one near these orbits
and supported in a tube containing no other quotient saddle.
It can be chosen one near every inner conditional minimum
over a smaller compact outer saddle neighborhood as well:
use continuity of the minimum maps, choose nested neighborhoods,
and average the cutoff over the compact based action.
The averaging leaves it one on the invariant inner neighborhood.
The quotient phase has a positive gap on the remaining compact
part of the support.

This cutoff is an auxiliary device for determining a leading
coefficient. It is neither the actual retained sharp window
nor a proposed global Wilson density.

\subsection{Exact Haar pushforward composition}
\label{olc-pushforward-section}

For $g>0$, consider the positive finite measure
\begin{equation}
              \chi(U)e^{-S_{\rm ext}(U)/g^2}\,d\mu_X(U).
 \label{olc-tube-measure}
\end{equation}
The maps $f,h,h\circ f$ are submersions near all relevant
orbits. At a small representative, the actual pivot blocks
give full rank; equivariance preserves rank on the orbit.
A smaller tube retains the property. Its local pushforward
densities are thus smooth relative to product Haar.

Write $J_g(V)$ for the pushforward density through $f$, and
$I_g(W)$ for the density through $h\circ f$. Exact
pushforward composition gives
\begin{equation}
                I_g(W)=\int_{h(V)=W}J_g(V)\,d\mu_h(V).
 \label{olc-Fubini}
\end{equation}
Here $d\mu_h$ is the conditional coarea density relative to
product Haar. Equivalently, test both sides against a smooth
function of $W$ and use ordinary Fubini. No multiplication
of formal delta distributions is required.

The first pushforward is exactly $H_L$-invariant. Indeed,
lift a based intermediate gauge by
\eqref{olc-group-sequence}, use the $H_{RL}$ invariance of
$\chi$ and $S_{\rm ext}$, and change the current Haar variables.
This is precisely the invariance needed in the outer integral.

The actual two-minimum comparisons identify the inner
critical value and then the outer one:
\begin{equation}
 \min_{f(U)=V}S_{\rm ext}(U)=S_{MR}(V),\qquad
 \min_{h(V)=W}S_{MR}(V)=S_{MRL}(W).
 \label{olc-minimum-values}
\end{equation}
These are local comparisons on the chosen tube.
The conditional minima depend smoothly on the outer
parameters modulo the inner based orbit. The cutoff is one
there. Outside the smaller outer saddle neighborhood, the
compact quotient phase gap makes the contribution
exponentially smaller. All uniformity here concerns the
fixed finite regulator and fixed ratios only.

\subsection{The finite-dimensional Laplace fact}
\label{olc-Laplace-section}

We record the elementary asymptotic used twice. In a finite
real chart suppose
\[
 F(z,w)=s(w)+\tfrac12z^TA(w)z+O(|z|^3),\qquad A(w)>0,
\]
and the critical point is the unique local minimum on the
chosen compact support. The density $m(z,w)$ is smooth and
its cutoff is one near zero. On a small compact parameter
germ the Hessian has a common positive lower bound. Then
\begin{equation}
 \int e^{-F(z,w)/g^2}m(z,w)\,dz
  =e^{-s(w)/g^2}g^n(2\pi)^{n/2}
        m(0,w)\det A(w)^{-1/2}[1+o(1)]
 \label{olc-Laplace}
\end{equation}
uniformly in that parameter germ.

To prove it, shrink the chart until $F-s\geq c|z|^2$.
Set $z=g\xi$, first restrict $\xi$ to a fixed bounded set,
and use Taylor's formula for uniform convergence to the
Gaussian. A common Gaussian majorant removes that restriction.
The compact remainder outside a smaller chart has a positive
phase gap and is exponentially negligible. This proof also
allows a smooth positive amplitude, whose leading value is
its value at the saddle, and an amplitude family converging
uniformly there with a bounded relative error.

For an orbit integral, first insert
\eqref{olc-gauge-insertion}. The phase becomes $S+G_K$.
Its unique saddle is the gauge-zero representative: the
quotient phase has its strict minimum there modulo the orbit,
and on the compact based orbit $G_K$ has its unique zero
at $P_K=I$. Away from a neighborhood of that representative,
the sum of the quotient and gauge penalties has a positive
gap. The actual retained chart therefore reduces the leading
asymptotic to \eqref{olc-Laplace}, with the exact factor
$c(g)^{-n_{s,K}}$ still present. No unspecified quotient
volume form or ultraviolet-uniform error estimate is used.

\subsection{Identification with the actual pivot and determinant}
\label{olc-identification-section}

The exact local constrained measure has, at its saddle,
the factor
\begin{equation}
                       h_0^{\,n_e-n_c}d_K(V_*).
 \label{olc-saddle-measure}
\end{equation}
Indeed relative Haar coordinates are
$d(e^{iB})=h_0j(B)\,dB$ with $j(0)=1$; the normalized
fluctuation straightening Jacobian is one at zero; and
solving the constraint contributes the inverse pivot product.
The target density is relative to group Haar. Consequently
there is no extra factor $j(\log W)$ in
\eqref{olc-saddle-measure}.

Let $I_{M,K;g}$ denote an appropriate local tube pushforward.
The phase-corrected, identity-normalized stationary amplitude is
\begin{equation}
 e^{D_{M,K}(W)}
   =\lim_{g\downarrow0}
       e^{[S_{MK}(W)-S_{MK}(1)]/g^2}
          \frac{I_{M,K;g}(W)}{I_{M,K;g}(1)} .
 \label{olc-normalized-amplitude}
\end{equation}
The gauge integral, powers of $g$, Haar constants, flat
determinants and all scalar cutoff factors cancel in this
ratio. The remaining saddle measure and Hessian are exactly
\eqref{olc-leading}. A cutoff equal to one near the critical
orbit changes no leading coefficient.

The actual sharp-window integral has the same saddle,
pivot factor and retained Hessian. Its variable is $z=g\xi$
and its window is the actual product window of threshold
$T(g)$. Choose $T(g)\to\infty$ with $gT(g)^3\to0$.
The established centered Gaussian estimate makes its
normalized logarithmic expectation
\begin{equation}
                          O\!\left(\tau(T(g))+g^2T(g)^6\right)
                              \longrightarrow0.
 \label{olc-sharp-leading}
\end{equation}
Thus the sharp-window and the auxiliary fixed-tube integrals
have the same normalized leading amplitude. This does not
equate the integrals at positive $g$.

The auxiliary first pushforward supplies only the intermediate
based $H_L$ invariance: its input was made $H_{RL}$-invariant.
The phase and normalized leading amplitude therefore have
that identity on the needed intermediate orbit. It suffices
to replace the rooted outer saddle by $V_p$.
The cutoff has not been made invariant under every output
site gauge with nontrivial final-center values. The full
local output gauge identity for the actual $D$ is the
previously proved one from the zero-input quantum leading
limit, not an additional invariance claimed for this cutoff.

\subsection{Sequential saddle evaluation}
\label{olc-proof-section}

\begin{proof}[Proof of Theorem~\ref{olc-theorem}]
Apply \eqref{olc-Laplace} on the $f$ fibers, uniformly over
a compact neighborhood of the outer stationary orbit.
With a scalar $c_1(g)$ independent of $V$,
\begin{equation}
 J_g(V)=c_1(g)e^{-S_{MR}(V)/g^2}
                   e^{D_{M,R}(V)}[1+o(1)] .
 \label{olc-inner-asymptotic}
\end{equation}
The relative error is uniform. Exact $H_L$ invariance reduces
the full compact outer orbit to a small slice. The inner
cutoff is one at every relevant conditional minimum.

Insert \eqref{olc-inner-asymptotic} into
\eqref{olc-Fubini}. The outer Laplace formula evaluates the
smooth amplitude $e^{D_{M,R}(V)}$ at its existing classical
minimum. With another scalar $c_2(g)$,
\begin{equation}
 I_g(W)=c_1(g)c_2(g)e^{-S_{MRL}(W)/g^2}
       e^{D_{M,R}(V_p^{MR,L}(W))+D_{MR,L}(W)}[1+o(1)] .
 \label{olc-sequential-asymptotic}
\end{equation}
The amplitude gradient has no factor $g^{-2}$ and hence
does not alter the leading phase saddle. The outer rooted
representative gives the same amplitude by the based
$H_L$ invariance. The compact phase gap controls the
remaining part of the outer integration.

Direct evaluation of the same integral through $h\circ f$
instead gives a scalar $c_{12}(g)$ and
\begin{equation}
 I_g(W)=c_{12}(g)e^{-S_{MRL}(W)/g^2}
                         e^{D_{M,RL}(W)}[1+o(1)] .
 \label{olc-direct-asymptotic}
\end{equation}
Normalize both equations by their values at $W=I$ after
removing the phase. There $V_p=I$ and every $D$ is zero.
All scalar factors cancel, without any computation of the
relation between $c_{12}$ and $c_1c_2$. Taking $g$ to zero
yields
\[
 e^{D_{M,RL}(W)}
   =e^{D_{M,R}(V_p^{MR,L}(W))+D_{MR,L}(W)} .
\]
On the positive real germ the logarithms are real and
zero at the identity. They therefore satisfy
\eqref{olc-cocycle}, without a branch ambiguity.
Both sides are given holomorphic pivot/determinant
functionals. The identity theorem extends their germ
identity to each connected common analytic domain.
The smaller comparison germ used by Laplace does not
shrink a domain already established for those functions.
\end{proof}

\subsection{Joint infinite systems and the Maxwell coefficient}
\label{olc-limit-section}

For fixed $R,L$, the existing modules give uniform anchored
analytic families for all three $D$ terms. On a joint local
limit system the finite identity passes on compact fields
by local normal convergence and anchored tail bounds;
the same bounds justify the classical substitutions.
It gives a scalar quotient cocycle, not a coefficientwise
identity of carrier decompositions.

Apply exact preservation of the true infinite Maxwell
coefficient under classical transport. Then
\begin{equation}
 \beta_{MRL}(D_{M,RL})
   =\beta_{MR}(D_{M,R})+\beta_{MRL}(D_{MR,L}).
 \label{olc-Maxwell-cocycle}
\end{equation}
Thus a long fixed-ratio leading coefficient is the sum of
its consecutive step coefficients in this joint system.
This proves neither their sign nor their convergence as
the preceding depth grows.

At finite volume, the symmetric short-cutoff moment
$\beta_{M,N}$ does not commute exactly with classical
transport. One must retain the separately controlled
finite leakage when applying it to
\eqref{olc-cocycle}. No exact finite scalar
$\beta$ additivity is asserted.

\subsection{Scope}
\label{olc-scope-section}

The result rests on a consistent local orbit-tube extension,
exact finite Haar pushforward composition, and ordinary
finite-dimensional saddle asymptotics using the actual
pivot, contour penalty and retained Hessian. It assumes
no ultraviolet-uniform stationary-phase remainder.

It does not identify direct and sequential sharp-window
sectors at positive $g$, extend the original local action
to every gauge configuration, or identify every local
sector with the full Wilson measure. The auxiliary cutoff
disappears only from the leading coefficient.

The background-field/ghost determinant comparison, a
bounded comparison error in depth, the numerical
logarithmic coefficient, its sign, physical coupling
tuning and the treatment of the other field sectors
remain separate from this cocycle.
